\documentclass[12pt,a4paper]{article}
\usepackage[spanish]{babel}
\usepackage[utf8]{inputenc}
\usepackage[T1]{fontenc}
\usepackage{amsmath, amssymb}
\usepackage{enumitem}
\usepackage{geometry}
\geometry{margin=2.5cm}

\title{Ejercicios 01\\
\large Modelos Y Arquitectura de Sistemas}
\author{Pitehr Hurtado-Cayo}

\begin{document}

\maketitle

\begin{enumerate}[label=\textbf{Pregunta \arabic*.}]
    \item \textbf{Se sabe que Descartes plantea la división del “todo”
              en sus partes, y que, por otro lado, sinergia es una
              propiedad emergente del “todo”. A partir de esto ¿diría Ud.
              que la concepción mecanicista al reconocer la existencia
              de un “todo” considera implícitamente la existencia de
              sinergia?}

          \textbf{Respuesta:} \\
          No, dado que la concepción mecanicista se enfoca en la división del "todo" en partes individuales,
          no reconoce explícitamente la existencia de sinergia. La sinergia es una propiedad emergente que
          surge de la interacción de las partes, y el mecanicismo tiende a analizar cada parte de manera aislada,
          sin considerar las interacciones complejas que pueden generar propiedades emergentes como la sinergia.

    \item \textbf{¿Es necesario que un objeto que puede ser
              considerado como un sistema, sea diseñado por personas
              con pensamiento sistémico?}

          \textbf{Respuesta:} \\
          No necesariamente, que un objeto puede ser considerado un sistema depende del observador, no de quien lo diseño.

    \item \textbf{Comente. “Muchos sistemas de interés son ‘sistemas abiertos’, por
              lo tanto, en realidad no es relevante reconocer límites en
              ellos, puesto siempre habrá entradas y salidas que
              atraviesen sus fronteras”.}

          \textbf{Respuesta:} \\
          FALSO. Distinguir los limites de un sistema es fundamental para distinguir al sistema de su entorno, por lo tanto, no es posible conceptualizar un sistema sin sus elementos, valga decir, limites.
    \item \textbf{Comente. ¿Puede un subsistema estar completamente aislado
              del resto de los subsistemas de un sistema?}

          \textbf{Respuesta:} \\
          NO. Un subsistema no puede estar completamente aislado del resto de los subsistemas de un sistema,
          ya que todos los subsistemas interactúan y dependen entre sí para el funcionamiento del sistema
          en su conjunto.

    \item \textbf{Comente. “Cuando se describe un sistema como sistema
              cerrado, se pone en riesgo la subsistencia del mismo,
              porque éste se volverá incapaz de incorporar las entradas
              que requiere”}

          \textbf{Respuesta:} \\
          FALSO. Primero, si al distinguir un sistema, donde distinguimos la necesidad de una entrada el entorno, no es un sistema cerrado.
          Segundo, al concepctualizar un sistema cerrado, no se pone en riesgo la subsistencia del mismo, sino del estudio en cuestion.


    \item \textbf{Considere la siguiente pregunta: ¿Es posible explicar
              sistémicamente el nacimiento y desarrollo de la ciencia
              moderna durante la revolución científica?
              A partir de la pregunta anterior seleccione el tipo de
              respuesta que considere correcto y desarróllelo:}
          \subitem \textbf{A. Sí es posible (Debe justificar su respuesta
              describiendo el nacimiento y desarrollo de la ciencia
              moderna desde una perspectiva sistémica).}
          \subitem \textbf{B. No es posible (Debe justificar su respuesta explicando
              por qué no es posible o por qué no tiene sentido la
              pregunta)}

          \textbf{Respuesta:} \\
          Si es posible, todos los fenomenos se pueden explicar desde una perspectiva sistémica, dado que el nacimiento y desarrollo de la ciencia moderna durante la revolución científica puede ser visto como un sistema complejo que involucra múltiples elementos interconectados, como los avances tecnológicos, las ideas filosóficas, las instituciones académicas y las interacciones sociales. Estos elementos interactúan entre sí de manera dinámica, dando lugar a la evolución y desarrollo de la ciencia moderna,
          partiendo desde la religion hasta la ciencia moderna.

    \item \textbf{Comente. “Si algo es un sistema, entonces no puede
              ser estudiado por un pensador reduccionista”.}

          \textbf{Respuesta:} \\
          FALSO, algo no puede SER un sistema, dado que los sistemas son conceptualizaciones del observador.
          Aunque un pensador reduccionista si podria estudiar un sistema a traves de la division del "todo" en partes, pero no podria comprender las propiedades emergentes del sistema, dado que estas surgen de la interacción de las partes y no pueden ser entendidas simplemente analizando cada parte por separado.

    \item \textbf{“Los objetos que tienen como característica
              la existencia de sinergia, deben ser estudiados
              desde la perspectiva sistémica”.}

          \textbf{Respuesta:} \\
          FALSO, la sinergia no es una caracteristica/atributo, sino una propiedad emergente.

    \item \textbf{“La existencia de sinergia es una caracterís2ca que se observa
              recursivamente en los sistemas”.}

          \textbf{Respuesta:} \\
          La sinergia no es una caracteristica/atributo, sino una propiedad emergente que se distingue, por lo tanto, el observador si puede
          distinguir sinergia recursivamente en los sistemas.

    \item \textbf{“La conducta teleológica que posee un sistema, puede considerarse
              como una propiedad sinérgica en él”}

          \textbf{Respuesta:} \\
          VERDADERO, la conducta teleológica que posee un sistema puede considerarse como una propiedad sinérgica en él, dado que la teleología se refiere a la capacidad de un sistema para tener un propósito o una función específica, y esta capacidad surge de la interacción de las partes del sistema, lo que es una característica de la sinergia.

    \item \textbf{El concepto “teleología” aplicado a un sistema nos indica que el
              “propósito del observador” es aquél elemento que no puede faltar en el
              estudio del mismo.}

          \textbf{Respuesta:} \\
          FALSO, el concepto de Teleologia se refiere al proposito o funcion de un sistema, no al proposito del observador.

    \item \textbf{“Todos los seres vivos finalmente mueren, por lo tanto, no se observan
              en la naturaleza sistemas realmente homeostácos”.}

          \textbf{Respuesta:} \\
          FALSO, La homeostasis se refiere a la capacidad de un sistema para mantener un equilibrio interno a pesar de las perturbaciones externas, no a la inmortalidad.
          Aunque todos los seres vivos eventualmente mueren, durante su vida pueden mantener un estado de homeostasis, lo que les permite sobrevivir y funcionar adecuadamente en su entorno.


    \item \textbf{¿Puede ocurrir que las variables fisiológicas de un sistema se mantengan
              en permanente variación, y sin embargo, éste se encuentre en estado
              estable?}

          \textbf{Respuesta:} \\
          VERDADERO. Existen los sistemas estables que oscilan.

    \item \textbf{Lea la sección 3.3 del libro “Introducción a la Teoría General de
              Sistemas” de Oscar Johansen. A partir de esta lectura, responda lo
              siguiente:}
          \subitem \textbf{A. Comente: “Puede ser difícil reconocer subsistemas, dado que no
              todas sus partes respetan el principio de recursividad. Por ejemplo,
              una uña no es un subsistema de un ser humano.” (p. 57)}

          \textbf{Respuesta:} \\
          FALSO. La recursividad no es un principio, sino un hecho que distingue el observador,
          Por lo tanto, segun Johansen si era necesario cumplir ciertas caracteristicas sistemicas,
          pero en TGS conceptualizar un subsistema no es una cuestion de ciertas caracteristicas, sino de la distincion del observador.

          \subitem \textbf{B. Comente: “Para reconocer subsistemas se deben identificar
              elementos presentes en un sistema. Un criterio para distinguir
              subsistemas, según Beer, es que estos sean sistemas viables”. (p.
              57)}

          \textbf{Respuesta:} \\
          FALSO. este criterio es útil cuando el observador analiza sistemas vivos u organizaciones; pero si el observador decide conceptualizar una estructura estática o un mecanismo de relojería
          , no utilizará la viabilidad como lente para distinguir sus subsistemas, sino que establecerá sus propios parámetros de acuerdo a su intención.

          \subitem \textbf{C. Comente: “Un hombre es un subsistema de una empresa, dado que
              él se encuentran los mismos subsistemas funcionales de Katz y Kahn
              que se identifican en las organizaciones” (p.58)}
          \textbf{Respuesta:} \\
          FALSO. Este enunciado es el ejemplo perfecto de cómo opera el observador. El hombre no es un subsistema de la empresa por una realidad ontológica o física indiscutible; lo es porque un observador intencionado decidió utilizar un modelo abstracto (el de Katz y Kahn) para conceptualizarlo de esa forma
          . Es el observador quien decide mapear las funciones biológicas del hombre (metabolismo, cerebro, etc.) y asociarlas conceptualmente a las áreas de la empresa (producción, dirección) para crear un paralelo recursivo
          . Si otro observador con un propósito distinto analizara la misma empresa —por ejemplo, un auditor financiero enfocándose solo en flujos de capital—, podría no considerar al hombre como un "subsistema recursivo", sino abstraerlo simplemente como un centro de costos o una variable numérica
          . La cualidad de subsistema se la otorga el modelo del observador.

    \item \textbf{De acuerdo con lo discutido en clases, seleccione la (o las) alternativa(s) que completa(n) correctamente la oración. En caso de que ninguna de las primeras 4 alternativas sea correcta, seleccione la alternativa “e”. Debe justificar su decisión. Si la alternativa seleccionada es la “e”, debe explicar por qué descarta cada una de las anteriores. La información no está sujeta a la “ley de conservación”, esto \underline{\hspace{3cm}}.}
          \begin{enumerate}
              \item implica que no debe representarse como entrada o salida en un diagrama de caja negra.
              \item ocurre porque cumple la ley de los incrementos.
              \item significa que la información no puede almacenarse.
              \item constituye una limitación del enfoque de sistemas.
              \item Ninguna de las anteriores completa correctamente la oración.
          \end{enumerate}
          \textbf{Respuesta:} \\
          De acuerdo con las fuentes proporcionadas, la alternativa correcta es la e) Ninguna de las anteriores completa correctamente la oración.
          A continuación, se presenta la justificación descartando cada una de las alternativas anteriores:
          \begin{itemize}
              \item Se descarta a) "implica que no debe representarse como entrada o salida en un diagrama de caja negra": Es una afirmación incorrecta. El diagrama de caja negra representa a un sistema interactuando con su entorno, y las fuentes establecen explícitamente que las corrientes de entrada y salida de un sistema abierto pueden ser perfectamente de materia, energía o información.
              \item Se descarta b) "ocurre porque cumple la ley de los incrementos": La formulación de esta alternativa establece una relación de causalidad falsa ("ocurre porque"). La información no está exenta de la ley de conservación debido a la ley de los incrementos. Por el contrario, la "ley de los incrementos" es simplemente el nombre descriptivo que el autor (Johansen) utiliza para explicar cómo se comporta empíricamente la información (donde hay una agregación neta y la salida no elimina la información del sistema). Las leyes describen fenómenos, no los causan. La información no responde a la ley de conservación física clásica sencillamente porque no es materia ni energía pura.
              \item Se descarta c) "significa que la información no puede almacenarse": Es una premisa falsa. Que no se comporte bajo la ley de conservación no impide su almacenamiento. De hecho, la mencionada "ley de los incrementos" afirma que en el sistema siempre existe una "agregación neta" entre la información que ya existe y la nueva que ingresa, demostrando que la información permanece y se acumula (almacena) en el sistema.
              \item Se descarta d) "constituye una limitación del enfoque de sistemas": Totalmente falso. La incapacidad de aplicar la ley de conservación a la información no limita al enfoque sistémico. Al contrario, la Teoría General de Sistemas y la Cibernética se nutren del concepto de información (a través de la comunicación, retroalimentación y la entropía negativa) para explicar cómo los organismos biológicos y las organizaciones sociales logran controlarse, adaptarse a su entorno y sobrevivir.
          \end{itemize}
    \item \textbf{
              En el libro Introducción a la Teoría General de Sistemas, se
              menciona que un sistema cerrado, según Forrester es aquel que
              cuenta con retroalimentación. ¿Quiere esto decir que los
              sistemas con retroalimentación están destinados a su
              desaparición por el principio de entropía?
              Explique cómo se manifiesta y cómo se puede combatir la
              entropía en una organización.
          }

          \textbf{Respuesta:} \\
          \textbf{Para responder a tus preguntas, es importante aclarar primero una confusión de conceptos que surge al utilizar el término ``sistema cerrado''.}
          \begin{enumerate}

              \item \textbf{¿Están los sistemas con retroalimentación destinados a desaparecer por la entropía?}

                    No, en absoluto. La aparente contradicción surge porque se están cruzando dos definiciones distintas de lo que es un ``sistema cerrado'':

                    \begin{itemize}
                        \item \textbf{El sistema cerrado según Forrester (circuito cerrado):} cuando autores como Forrester hablan de un ``sistema cerrado'', en realidad tienen en mente un sistema con ``circuito cerrado''. Un circuito cerrado es precisamente el mecanismo de retroalimentación (o \textit{feedback}), el cual permite al sistema autorregularse y compensar desviaciones para alcanzar sus objetivos.

                        \item \textbf{El sistema cerrado termodinámico:} por otro lado, la física define a un sistema cerrado como aquel que no intercambia materia, energía ni información con su entorno. Es a este tipo de sistemas al que se le aplica de forma estricta la Segunda Ley de la Termodinámica, la cual dicta que su entropía (desorden) aumentará irremediablemente hasta llegar a su fin.
                    \end{itemize}

                    Las organizaciones (y los seres vivos) son en realidad sistemas abiertos, ya que necesitan importar constantemente recursos de su medio para sobrevivir. Lejos de estar destinados a desaparecer, utilizar el mecanismo de retroalimentación (el ``circuito cerrado'' de Forrester) es la herramienta fundamental que les permite ajustar su comportamiento, mantener su estabilidad y evitar su destrucción.

              \item \textbf{¿Cómo se manifiesta la entropía en una organización?}

                    En términos generales, la entropía se manifiesta como la tendencia hacia el máximo desorden, la desorganización y el caos.

                    Llevado al ámbito organizacional y de la comunicación, la entropía se manifiesta operativamente como incertidumbre (falta de información). Una organización sufre de entropía cuando se enfrenta a una alta incertidumbre a la hora de tomar decisiones críticas (por ejemplo: ¿debe invertir?, ¿debe cambiar de proveedor?, ¿debe lanzar un nuevo producto?). A mayor incertidumbre, mayor es la entropía y la desorganización interna.

              \item \textbf{¿Cómo se puede combatir la entropía en la organización?}

                    Para que una organización o sistema vivo no sucumba ante la entropía, debe importar desde su entorno lo que se conoce como neguentropía (entropía negativa). Esto se logra principalmente a través de dos elementos:

                    \begin{itemize}
                        \item \textbf{Importación de energía:} para mantenerse y no desaparecer, la organización necesita incorporar energía y recursos de su entorno que le permitan sostener sus procesos.

                        \item \textbf{Adquisición de información:} la incertidumbre en una organización se combate con información. La Teoría de la Información establece que la información es matemáticamente isomórfica (equivalente) a la neguentropía; es decir, la información es una medida de organización. Por lo tanto, la información actúa como un ``agente neguentrópico'' directo: mientras más información reciba y procese la organización sobre su entorno y sus propios resultados, menor será su incertidumbre y mayor será su capacidad para mantenerse organizada y estable.
                    \end{itemize}

          \end{enumerate}

    \item   \textbf{Puede el pensamiento científico-reduccionista del mecanicismo haber influido en el
              nacimiento del paradigma sistémico? Explique}

          \textbf{Respuesta:} \\
          Sí, el pensamiento científico-reduccionista del mecanicismo influyó profundamente en el nacimiento del paradigma sistémico, pero lo hizo como una respuesta y reacción frente a sus severas limitaciones.
          La influencia se dio de las siguientes maneras:

          \textbf{La incapacidad de explicar la totalidad y la sinergia:} El paradigma mecanicista se basaba en el enfoque analítico y el reduccionismo; es decir, intentaba comprender el universo fragmentándolo en sus unidades más pequeñas y elementales, analizándolas de forma aislada. Sin embargo, a medida que la ciencia comenzó a estudiar fenómenos más complejos (como los organismos biológicos, la psicología o la sociedad), este método resultó ser insuficiente. Al aislar las partes, el mecanicismo destruía las propiedades emergentes o sinérgicas que nacen exclusivamente de la interacción dinámica de los elementos. El paradigma sistémico nació ante la estricta necesidad de abordar estos problemas de ``totalidad'' y de organización que el reduccionismo dejaba fuera.
          \textbf{La respuesta a un ``mundo muerto'' sin propósito:} El mecanicismo clásico, inspirado en la física newtoniana y el espíritu laplaciano, presentaba al universo como un juego de partículas gobernadas por leyes causales ciegas e inexorables. Esta visión dejaba por fuera las cualidades, la consciencia, el espíritu y, sobre todo, el propósito o direccionalidad (telos). El paradigma sistémico reaccionó reintroduciendo la teleología (comportamiento orientado a metas) y reconociendo a los organismos como sistemas abiertos, alejándose del rígido determinismo mecánico.
          \textbf{La crisis de la especialización y la incomunicación:} El reduccionismo llevó a que la ciencia se dividiera en innumerables disciplinas y subdisciplinas, encapsulando a los científicos en universos privados donde cada uno hablaba su propio lenguaje. Esto generó una ``crisis del conocimiento'' provocada por una ``sordera especializada'', amenazando con desintegrar la ciencia. La Teoría General de Sistemas nació motivada por este problema, proponiéndose como un ``oído generalizado''. Su objetivo fue buscar isomorfismos (modelos, conceptos y leyes comunes aplicables a diferentes campos) para restablecer la comunicación, evitar la repetición de esfuerzos y devolverle la unidad a la ciencia.

          En resumen, el mecanicismo empujó a la ciencia hasta un límite donde fragmentar más las cosas ya no aportaba respuestas útiles para la complejidad de la realidad, obligando al surgimiento de una nueva visión (la sistémica) que pudiera integrar y observar las cosas como totalidades.

\end{enumerate}

\end{document}